\clearpage
\section*{Final comparison}
\begin{table}[H]\centering\normalsize
\caption{Ten-model validation comparison. Layout specifies tree count and
maximum depth; D4 means depth four. $P_{\rm nz}$ counts learned nonzero scalar
values. $\epsilon_s$ is interpolated at 1\% background;
$\epsilon_s^{\rm strict}$ uses a strict threshold accepting at most 525
background events. C++ times are measured with the same sparse reference
kernel. $M_{\rm path}$ is a selected-path upper bound and $M_{\rm all}$ counts
all structurally reachable affine functions, both after exact-zero and
$\pm1$ factor elimination.}
\label{tab:comparison}
\setlength{\tabcolsep}{4pt}\renewcommand{\arraystretch}{1.28}
\begin{tabular}{@{}llrrrrrrr@{}}\toprule
Family & Layout & $P_{\rm nz}$ & \shortstack{$\epsilon_s$\\{[\%]}} & \shortstack{$\epsilon_s^{\rm strict}$\\{[\%]}} & AUC & \shortstack{C++\\{[$\mu$s/event]}} & $M_{\rm path}$ & $M_{\rm all}$\\\midrule
XGBoost & 100$\times$D4 & 3099 & 55.48 & 55.48 & 0.92056 & 3.675 & 0 & 0 \\
\midrule
Axis / linear & 3$\times$D4 & 479 & 56.32 & 56.32 & 0.92576 & 0.101 & 35 & 386 \\
Axis / linear & 1$\times$D4 & 228 & 54.74 & 54.74 & 0.92387 & 0.048 & 16 & 197 \\
Axis / linear & 1$\times$D3 & 132 & 50.10 & 50.10 & 0.92020 & 0.049 & 16 & 117 \\
\midrule
Oblique / constant & 5$\times$D4 & 1250 & 55.40 & 54.75 & 0.92610 & 0.965 & 272 & 1088 \\
Oblique / constant & 3$\times$D4 & 753 & 53.54 & 53.30 & 0.92318 & 0.542 & 160 & 656 \\
Oblique / constant & 1$\times$D4 & 256 & 45.36 & 45.33 & 0.91247 & 0.171 & 64 & 224 \\
\midrule
Oblique / linear & 5$\times$D4 & 1275 & 57.73 & 57.73 & 0.93090 & 0.412 & 212 & 1096 \\
Oblique / linear & 3$\times$D4 & 837 & 57.12 & 57.12 & 0.92967 & 0.320 & 160 & 734 \\
Oblique / linear$^{c}$ & 3$\times$D4 & 418 & 56.20 & 56.20 & 0.92733 & 0.167 & 78 & 325 \\
\bottomrule

\end{tabular}

\vspace{8pt}
\begin{minipage}{\textwidth}\small
$^{c}$Compact refit. Axis rows additionally use 45, 15 and 7 feature indices,
respectively, and one direction bit per nonconstant cut. Their fixed unit
routing coefficients are excluded from $P_{\rm nz}$. XGBoost stores 1499
feature indices; sparse oblique storage can also require coordinate indices.
These scalar counts are not complete memory-footprint estimates.

All selected TAO snapshots come from converged source stages; selection can choose
an earlier checkpoint. The axis 3$\times$D4 entry uses its regularized source:
its unpenalized leaf refit was held for separation. The 837-coefficient
oblique/linear entry uses a converged leaf-only stage; its held full stage is
excluded. Representatives cover different scalar budgets within each family.

Timing: one CPU thread, 132,800 events, median of nine repeats after warmup;
parsing, input conversion and I/O excluded. Native optimized XGBoost separately
measured 1.478 $\mu$s/event, including probability output. Working points are calibrated to the stated
reference arithmetic and must be rechecked after a change of reduction order
or precision. The common C++
column is not a ranking of maximally optimized software implementations.

Multiplication counts exclude preprocessing, sigmoid, additions and
comparisons; no shift or general constant-multiplier optimization is assumed
beyond exact-zero and unit factors. They are arithmetic workloads, not
synthesized DSP/LUT/BRAM use or FPGA latency. Results use the fixed validation
split and do not include a new test evaluation.
\end{minipage}
\end{table}
